\chapter{Complexity — Logspace Uniform Walk}
%%% generated by the blueprint pipeline from TCSlib.Complexity.CircuitComplexity.LogspaceUniformWalk
\section{Overview}

Declarations of the Lean module \texttt{TCSlib.Complexity.CircuitComplexity.LogspaceUniformWalk}.

\section{Declarations}

\begin{definition}[Adjacency-walk sites]
\lean{BoolCircuit.AdjWalk.Site}
\label{BoolCircuit.AdjWalk.Site}
\leanok
Define the set of sites to consist of the sixteen positions named inOne, inEnd, gOne,
gEnd, kA1, kA2, kO1, kO2, kN1, kN2, aOne, uOne, uEnd, aEnd, oOne, and oEnd.
\end{definition}

\begin{definition}[Adjacency-walk labels]
\lean{BoolCircuit.AdjWalk.Lbl}
\label{BoolCircuit.AdjWalk.Lbl}
\leanok
\uses{BoolCircuit.AdjWalk.Site}
Define the label set to contain the labels start, nI, incV, gS, kA, kB, aS, aL, aE,
clrW, uL, incW, aN, incV2, oS, oI, oL, incWo, and done; in addition, it contains em and
bump labels indexed by a site, and hitL labels indexed by a Boolean value.
\end{definition}

\begin{definition}[Bit assigned to each site]
\lean{BoolCircuit.AdjWalk.Site.bit}
\label{BoolCircuit.AdjWalk.Site.bit}
\leanok
\uses{BoolCircuit.AdjWalk.Site}
Define a Boolean value for each site that is true at inOne, gOne, kO2, kN1, aOne, uOne,
and oOne, and false at every other site.
\end{definition}

\begin{definition}[Successor label of a site]
\lean{BoolCircuit.AdjWalk.Site.next}
\label{BoolCircuit.AdjWalk.Site.next}
\leanok
\uses{BoolCircuit.AdjWalk.Lbl, BoolCircuit.AdjWalk.Site}
The successor label assigns to the sites, respectively: increment the vertex register at
the first input site; begin the gate scan at the final input site; test the gate kind at
the first gate site; begin the output scan at the final gate site; emit the second site
of each of the three kind-test pairs at the first site of that pair; begin the argument
scan at the second site of each pair; clear the work register at the first argument
site; increment the work register at the first unary-input site; begin the next argument
at the final unary-input site; increment the vertex register at the final argument site;
increment the output-scan register at the first output site; and halt at the final
output site.
\end{definition}

\begin{definition}[Adjacency-walk input languages]
\lean{BoolCircuit.AdjWalk.decs}
\label{BoolCircuit.AdjWalk.decs}
\leanok
\uses{BoolCircuit.DAGCircuitFamily, BoolCircuit.DAGCircuitFamily.edgeLang, BoolCircuit.DAGCircuitFamily.sizeLang, BoolCircuit.DAGCircuitFamily.typeLang}
For a family of Boolean circuits, the five languages indexed by $0,1,2,3,4$ are,
respectively, its size language, the type language for input vertices, the type language
for AND gates, the type language for OR gates, and its edge language.
\end{definition}

\begin{definition}[Adjacency-walk machine]
\lean{BoolCircuit.AdjWalk.walk}
\label{BoolCircuit.AdjWalk.walk}
\leanok
\uses{BoolCircuit.AdjWalk.Lbl, BoolCircuit.AdjWalk.Site.bit, BoolCircuit.AdjWalk.Site.next, Complexity.LogProg.ARM}
For a Boolean parameter $\ell$, define a four-register machine with five available
deciders and the labels in the walk control graph as follows. At each label it performs
the specified register operation, branch, decider call, or return: it validates the
first input payload; scans input and gate indices; tests gate kinds and arguments; scans
circuit edges and output vertices; and, at an emitted site, compares the input payload
with register zero, returning the disjunction of $\ell$ and the site's bit on a match,
or incrementing register zero and proceeding to the site's successor label otherwise.
\end{definition}

\begin{definition}[Four-coordinate register parameters]
\lean{BoolCircuit.AdjWalk.rg}
\label{BoolCircuit.AdjWalk.rg}
\leanok
For natural numbers $k,V,U,W$, define $r:\text{Fin }4\to\bbn$ by $r(0)=k$, $r(1)=V$,
$r(2)=U$, and $r(3)=W$.
\end{definition}

\begin{lemma}[First register coordinate]
\lean{BoolCircuit.AdjWalk.rg_0}
\label{BoolCircuit.AdjWalk.rg_0}
\leanok
\uses{BoolCircuit.AdjWalk.rg}
\difficulty{1}
For natural numbers $k,V,U,W$, the first coordinate of the four-register tuple
$(k,V,U,W)$ is $k$.
\end{lemma}

\begin{lemma}[Second register coordinate]
\lean{BoolCircuit.AdjWalk.rg_1}
\label{BoolCircuit.AdjWalk.rg_1}
\leanok
\uses{BoolCircuit.AdjWalk.rg}
\difficulty{1}
For natural numbers $k,V,U,W$, the second coordinate of the four-register tuple
$(k,V,U,W)$ is $V$.
\end{lemma}

\begin{lemma}[Third register coordinate]
\lean{BoolCircuit.AdjWalk.rg_2}
\label{BoolCircuit.AdjWalk.rg_2}
\leanok
\uses{BoolCircuit.AdjWalk.rg}
\difficulty{1}
For natural numbers $k,V,U,W$, the third coordinate of the four-register tuple
$(k,V,U,W)$ is $U$.
\end{lemma}

\begin{lemma}[Fourth register coordinate]
\lean{BoolCircuit.AdjWalk.rg_3}
\label{BoolCircuit.AdjWalk.rg_3}
\leanok
\uses{BoolCircuit.AdjWalk.rg}
\difficulty{1}
For natural numbers $k,V,U,W$, the fourth coordinate of the four-register tuple
$(k,V,U,W)$ is $W$.
\end{lemma}

\begin{lemma}[Exhaustion of four registers]
\lean{BoolCircuit.AdjWalk.fin4_cases}
\label{BoolCircuit.AdjWalk.fin4_cases}
\leanok
\difficulty{4}
If a property holds for each of the four registers numbered $0,1,2,3$, then it holds for
every register among the four.
\end{lemma}

\begin{lemma}[Updating the first coordinate]
\lean{BoolCircuit.AdjWalk.upd_0}
\label{BoolCircuit.AdjWalk.upd_0}
\leanok
\uses{BoolCircuit.AdjWalk.fin4_cases, BoolCircuit.AdjWalk.rg}
\difficulty{4}
For natural numbers $k,V,U,W,q$, let $r(k,V,U,W)$ be the function on $\text{Fin}(4)$ with
successive values $k,V,U,W$. Replacing its value at index $0$ by $q$ gives $r(q,V,U,W)$.
\end{lemma}

\begin{lemma}[Updating the second coordinate]
\lean{BoolCircuit.AdjWalk.upd_1}
\label{BoolCircuit.AdjWalk.upd_1}
\leanok
\uses{BoolCircuit.AdjWalk.fin4_cases, BoolCircuit.AdjWalk.rg}
\difficulty{4}
For natural numbers $k,V,U,W,q$, let $r(k,V,U,W)$ be the function on $\text{Fin}(4)$ with
successive values $k,V,U,W$. Replacing its value at index $1$ by $q$ gives $r(k,q,U,W)$.
\end{lemma}

\begin{lemma}[Updating the third coordinate]
\lean{BoolCircuit.AdjWalk.upd_2}
\label{BoolCircuit.AdjWalk.upd_2}
\leanok
\uses{BoolCircuit.AdjWalk.fin4_cases, BoolCircuit.AdjWalk.rg}
\difficulty{4}
For natural numbers $k,V,U,W,q$, let $r(k,V,U,W)$ be the function on $\text{Fin}(4)$ with
successive values $k,V,U,W$. Replacing its value at index $2$ by $q$ gives $r(k,V,q,W)$.
\end{lemma}

\begin{lemma}[Updating the fourth coordinate]
\lean{BoolCircuit.AdjWalk.upd_3}
\label{BoolCircuit.AdjWalk.upd_3}
\leanok
\uses{BoolCircuit.AdjWalk.fin4_cases, BoolCircuit.AdjWalk.rg}
\difficulty{4}
For natural numbers $k,V,U,W,q$, let $r(k,V,U,W)$ be the function on $\text{Fin}(4)$ with
successive values $k,V,U,W$. Replacing its value at index $3$ by $q$ gives $r(k,V,U,q)$.
\end{lemma}

\begin{lemma}[Initial register tuple]
\lean{BoolCircuit.AdjWalk.rg_init}
\label{BoolCircuit.AdjWalk.rg_init}
\leanok
\uses{BoolCircuit.AdjWalk.fin4_cases, BoolCircuit.AdjWalk.rg}
\difficulty{4}
The constant function on the four register indices with value zero equals the register
tuple $(0,0,0,0)$.
\end{lemma}

\begin{definition}[Configuration register bound]
\lean{BoolCircuit.AdjWalk.Gd}
\label{BoolCircuit.AdjWalk.Gd}
\leanok
\uses{BoolCircuit.AdjWalk.Lbl, Complexity.LogProg.AConf}
For a natural number $B$ and a configuration of the four-register walk machine, the
configuration is bounded by $B$ if every register value is at most $B$.
\end{definition}

\begin{lemma}[Register Bounds for the Register Tuple]
\lean{BoolCircuit.AdjWalk.gd_rg}
\label{BoolCircuit.AdjWalk.gd_rg}
\leanok
\uses{BoolCircuit.AdjWalk.Gd, BoolCircuit.AdjWalk.Lbl, BoolCircuit.AdjWalk.fin4_cases, BoolCircuit.AdjWalk.rg}
\difficulty{4}
Let $B,k,V,U,W$ be natural numbers, and let $\ell$ be an optional control label and $r$
an optional Boolean result. If $k,V,U,W\le B$, then the configuration with label $\ell$,
registers $(k,V,U,W)$, and result $r$ has every register at most $B$.
\end{lemma}

\begin{definition}[Encoded unary and binary input]
\lean{BoolCircuit.AdjWalk.X}
\label{BoolCircuit.AdjWalk.X}
\leanok
\uses{Turing.pairEncode}
For natural numbers $n$ and $i$, write $X$ for the pairing of the unary word of length
$n$ and the canonical binary word for $i$, namely $\langle
1^n,\operatorname{bits}(i)\rangle$.
\end{definition}

\begin{definition}[Emitted-site progress condition]
\lean{BoolCircuit.AdjWalk.Em}
\label{BoolCircuit.AdjWalk.Em}
\leanok
\uses{BoolCircuit.AdjWalk.Gd, BoolCircuit.AdjWalk.Lbl, BoolCircuit.AdjWalk.X, BoolCircuit.AdjWalk.walk, Complexity.LogProg.AConf, Complexity.LogProg.AHalt, Complexity.LogProg.AReach}
Fix a Boolean parameter $\ell$, natural numbers $n,i,B$, a Boolean-valued function $o$
on five indices and bit strings, a list $S$ of Boolean values, a configuration $a$, and
a predicate $P$ on configurations. The condition $\mathrm{Em}(k,S,a,P)$ requires both of
the following: whenever $k\le i<k+|S|$, the walk machine halts from $a$ on input
$\langle 1^n,\operatorname{bits}(i)\rangle$ with outcome $\ell\lor S[i-k]$ (using false
if that list position is out of range), under the register bound $B$; and whenever
$k+|S|\le i$, there is a configuration satisfying $P$ that is reachable from $a$ on that
input under the same register bound.
\end{definition}

\begin{lemma}[Concatenation of emitted segments]
\lean{BoolCircuit.AdjWalk.Em.append}
\label{BoolCircuit.AdjWalk.Em.append}
\leanok
\uses{BoolCircuit.AdjWalk.Em, BoolCircuit.AdjWalk.Lbl, Complexity.LogProg.AConf, Complexity.LogProg.AReach.halt, Complexity.LogProg.AReach.trans}
\difficulty{5}
Fix a Boolean parameter $\ell$, natural numbers $n,i,B$, and an oracle assigning a
Boolean answer to each of five deciders and each bit string. For a segment $S$ starting
at position $k$, let $a$ satisfy the emission property with postcondition $P$; suppose
also that whenever $k+|S|\le i$, every configuration satisfying $P$ satisfies the
emission property for a segment $T$ starting at $k+|S|$ with postcondition $Q$. Then the
concatenated segment $S T$ starting at $k$ satisfies the emission property from $a$ with
postcondition $Q$.
\end{lemma}

\begin{lemma}[Empty segment emission property]
\lean{BoolCircuit.AdjWalk.Em.nil}
\label{BoolCircuit.AdjWalk.Em.nil}
\leanok
\uses{BoolCircuit.AdjWalk.Em, BoolCircuit.AdjWalk.Lbl, Complexity.LogProg.AConf, Complexity.LogProg.AReach.refl}
\difficulty{2}
Fix a Boolean parameter $\ell$, natural numbers $n,i,B,k$, and an oracle assigning a
Boolean answer to each of five deciders and each bit string. If the initial
configuration $a$ satisfies $P$, then the empty segment starting at $k$ satisfies the
emission property from $a$ with postcondition $P$.
\end{lemma}

\begin{lemma}[Prepending a bounded reach]
\lean{BoolCircuit.AdjWalk.Em.pre}
\label{BoolCircuit.AdjWalk.Em.pre}
\leanok
\uses{BoolCircuit.AdjWalk.Em, BoolCircuit.AdjWalk.Gd, BoolCircuit.AdjWalk.Lbl, BoolCircuit.AdjWalk.X, BoolCircuit.AdjWalk.walk, Complexity.LogProg.AConf, Complexity.LogProg.AReach, Complexity.LogProg.AReach.halt, Complexity.LogProg.AReach.trans}
\difficulty{4}
Fix a Boolean parameter $\ell$, natural numbers $n,i,B$, and an oracle assigning a
Boolean answer to each of five deciders and each bit string. If $a$ can reach $b$ while
maintaining the register bound, and $b$ satisfies the emission property for a segment
$S$ starting at $k$ with postcondition $P$, then $a$ satisfies that same emission
property.
\end{lemma}

\begin{lemma}[One-step predecessor for emission]
\lean{BoolCircuit.AdjWalk.Em.step}
\label{BoolCircuit.AdjWalk.Em.step}
\leanok
\uses{BoolCircuit.AdjWalk.Em, BoolCircuit.AdjWalk.Em.pre, BoolCircuit.AdjWalk.Gd, BoolCircuit.AdjWalk.Lbl, BoolCircuit.AdjWalk.X, BoolCircuit.AdjWalk.walk, Complexity.LogProg.AConf, Complexity.LogProg.AReach.single, Complexity.LogProg.astep}
\difficulty{1}
Fix a Boolean parameter $\ell$, natural numbers $n,i,B$, and an oracle assigning a
Boolean answer to each of five deciders and each bit string. If the registers of a
configuration $a$ are all at most $B$, and the configuration obtained by one step of the
specified walk from $a$ satisfies the emission property for a segment $S$ starting at
$k$ with postcondition $P$, then $a$ itself satisfies that emission property.
\end{lemma}

\begin{lemma}[Monotonicity of the emission property]
\lean{BoolCircuit.AdjWalk.Em.mono}
\label{BoolCircuit.AdjWalk.Em.mono}
\leanok
\uses{BoolCircuit.AdjWalk.Em, BoolCircuit.AdjWalk.Lbl, Complexity.LogProg.AConf}
\difficulty{4}
Fix a Boolean parameter $\ell$, natural numbers $n,i,B$, and an oracle assigning a
Boolean answer to each of five deciders and each bit string. If a configuration
satisfies the emission property for a segment $S$ starting at $k$ with postcondition
$P$, and every configuration satisfying $P$ also satisfies $Q$, then it satisfies the
emission property for the same segment and starting position with postcondition $Q$.
\end{lemma}

\begin{lemma}[Continuation of an emission property]
\lean{BoolCircuit.AdjWalk.Em.cont}
\label{BoolCircuit.AdjWalk.Em.cont}
\leanok
\uses{BoolCircuit.AdjWalk.Em, BoolCircuit.AdjWalk.Gd, BoolCircuit.AdjWalk.Lbl, BoolCircuit.AdjWalk.X, BoolCircuit.AdjWalk.walk, Complexity.LogProg.AConf, Complexity.LogProg.AReach, Complexity.LogProg.AReach.trans}
\difficulty{4}
Fix a Boolean parameter $\ell$, natural numbers $n,i,B$, and an oracle assigning a
Boolean answer to each of five deciders and each bit string. Suppose $a$ satisfies the
emission property for a segment $S$ starting at $k$ with postcondition $P$, and suppose
that whenever $k+|S|\le i$, every configuration satisfying $P$ can reach a configuration
satisfying $Q$ while maintaining the register bound. Then $a$ satisfies the emission
property for the same segment $S$ and starting position $k$, now with postcondition $Q$.
\end{lemma}

\begin{lemma}[Emission of a single site bit]
\lean{BoolCircuit.AdjWalk.Em.emit}
\label{BoolCircuit.AdjWalk.Em.emit}
\leanok
\uses{BoolCircuit.AdjWalk.Em, BoolCircuit.AdjWalk.Site, BoolCircuit.AdjWalk.Site.bit, BoolCircuit.AdjWalk.Site.next, BoolCircuit.AdjWalk.X, BoolCircuit.AdjWalk.gd_rg, BoolCircuit.AdjWalk.rg, BoolCircuit.AdjWalk.rg_0, BoolCircuit.AdjWalk.upd_0, BoolCircuit.AdjWalk.walk, Complexity.LogProg.AHalt.ret, Complexity.LogProg.AHalt.step, Complexity.LogProg.AReach.refl, Complexity.LogProg.AReach.step, Complexity.LogProg.astep, Complexity.LogProg.bits_injective, Complexity.LogProg.plainWord_pairEncode}
\difficulty{5}
Fix a Boolean parameter $\ell$, natural numbers $n,i,B$, an oracle assigning a Boolean
answer to each of five deciders and each bit string, and a site $s$. Let $k,V,U,W$ be
natural numbers with $k\le i$, $k+1\le B$, and $V,U,W\le B$. Starting from the
configuration with control state $\mathrm{em}(s)$, registers $(k,V,U,W)$, and no answer,
the one-term segment consisting of the bit assigned to $s$ satisfies the emission
property, with postcondition that the resulting configuration has control state $s$'s
successor, registers $(k+1,V,U,W)$, and no answer.
\end{lemma}

\begin{definition}[Circuit decider family]
\lean{BoolCircuit.AdjWalk.orc}
\label{BoolCircuit.AdjWalk.orc}
\leanok
\uses{BoolCircuit.AdjWalk.decs, BoolCircuit.DAGCircuitFamily, Turing.MultiTapeTM.indicator}
For a family $C$ of Boolean circuits, define its Boolean outcome function on an index
$j$ among five choices and a bit string $V$ to be true exactly when $V$ belongs to the
$j$th language associated with $C$.
\end{definition}

\begin{lemma}[One-argument oracle-call step]
\lean{BoolCircuit.AdjWalk.astep_call₁}
\leanok
\uses{BoolCircuit.AdjWalk.Lbl, BoolCircuit.AdjWalk.X, BoolCircuit.AdjWalk.walk, Complexity.LogProg.astep, Complexity.LogProg.vword_unary₁, Turing.pairEncode}
\difficulty{3}
Let $\ell$ be a label, $j\in\{0,1,2,3,4\}$, $r\in\{0,1,2,3\}$, and let $v$ assign a
natural number to each of the four registers. If the instruction at $\ell$ calls oracle
$j$ in unary-first mode with the single register $r$, with successor labels $\ell_1$ and
$\ell_0$ for positive and negative answers, respectively, then one machine step from the
configuration with label $\ell$, registers $v$, and no answer leaves the registers
unchanged and no answer recorded, and moves to $\ell_1$ exactly when the oracle accepts
the encoding of the unary string $1^n$ paired with the binary representation of $v(r)$;
otherwise it moves to $\ell_0$.
\end{lemma}

\begin{lemma}[Two-argument oracle-call step]
\lean{BoolCircuit.AdjWalk.astep_call₂}
\leanok
\uses{BoolCircuit.AdjWalk.Lbl, BoolCircuit.AdjWalk.X, BoolCircuit.AdjWalk.walk, Complexity.LogProg.astep, Complexity.LogProg.vword_unary₂, Turing.pairEncode}
\difficulty{3}
Let $\ell$ be a label, $j\in\{0,1,2,3,4\}$, and $r,r'\in\{0,1,2,3\}$, and let $v$ assign
a natural number to each of the four registers. If the instruction at $\ell$ calls
oracle $j$ in unary-first mode with the ordered registers $r,r'$, with successor labels
$\ell_1$ and $\ell_0$ for positive and negative answers, respectively, then one machine
step from the configuration with label $\ell$, registers $v$, and no answer leaves the
registers unchanged and no answer recorded, and moves to $\ell_1$ exactly when the
oracle accepts the encoding of $1^n$ paired with the encoding of the binary
representations of $v(r)$ and $v(r')$; otherwise it moves to $\ell_0$.
\end{lemma}

\begin{lemma}[Oracle Answer for Circuit Size]
\lean{BoolCircuit.AdjWalk.orc_size}
\label{BoolCircuit.AdjWalk.orc_size}
\leanok
\uses{BoolCircuit.AdjWalk.decs, BoolCircuit.AdjWalk.orc, BoolCircuit.DAGCircuit.size, BoolCircuit.DAGCircuitFamily.mem_sizeLang, Complexity.LogProg.indicator_eq_decide, Turing.pairEncode}
\difficulty{3}
Let $C$ be a family of Boolean circuits with $n$ input vertices, and let $v$ be a
natural number. The size oracle, queried with the encoding of $n$ and the binary
encoding of $v$, returns true exactly when $v$ is less than the size of the circuit for
input length $n$.
\end{lemma}

\begin{lemma}[Type query for encoded vertices]
\lean{BoolCircuit.AdjWalk.orc_type}
\label{BoolCircuit.AdjWalk.orc_type}
\leanok
\uses{BoolCircuit.AdjWalk.decs, BoolCircuit.AdjWalk.orc, BoolCircuit.DAGCircuit.size, BoolCircuit.DAGCircuit.vtype, BoolCircuit.DAGCircuitFamily.mem_typeLang, BoolCircuit.DAGCircuitFamily.typeLang, BoolCircuit.GateKind, Complexity.LogProg.indicator_eq_decide, Turing.pairEncode}
\difficulty{3}
Let $C$ be a family of directed acyclic circuits, let $j$ be one of the five query
indices, let $t$ be an optional gate kind, and let $n,v$ be natural numbers. If the
query language at index $j$ is the language of vertices of type $t$, then the Boolean
indicator of that query language on the encoding of $1^n$ paired with the binary
expansion of $v$ equals true exactly when $v$ is less than the size of the circuit on
$n$ inputs and vertex $v$ has type $t$.
\end{lemma}

\begin{lemma}[Oracle Answer for Circuit Edges]
\lean{BoolCircuit.AdjWalk.orc_edge}
\label{BoolCircuit.AdjWalk.orc_edge}
\leanok
\uses{BoolCircuit.AdjWalk.decs, BoolCircuit.AdjWalk.orc, BoolCircuit.DAGCircuit.Edge, BoolCircuit.DAGCircuitFamily.mem_edgeLang, Complexity.LogProg.indicator_eq_decide, Turing.pairEncode}
\difficulty{3}
Let $C$ be a family of Boolean circuits with $n$ input vertices, and let $u,v$ be
natural numbers. The edge oracle, queried with the encoding of $n$ and the pair of
binary encodings of $u$ and $v$, returns true exactly when $u$ is an argument of the
gate at vertex $v$ in the circuit for input length $n$; equivalently, its Boolean answer
is the indicator of that edge relation.
\end{lemma}

\begin{lemma}[Walk behavior in the unary loop]
\lean{BoolCircuit.AdjWalk.uLoop}
\label{BoolCircuit.AdjWalk.uLoop}
\leanok
\uses{BoolCircuit.AdjWalk.Em, BoolCircuit.AdjWalk.Em.append, BoolCircuit.AdjWalk.Em.emit, BoolCircuit.AdjWalk.Em.mono, BoolCircuit.AdjWalk.Em.step, BoolCircuit.AdjWalk.Site.bit, BoolCircuit.AdjWalk.Site.next, BoolCircuit.AdjWalk.gd_rg, BoolCircuit.AdjWalk.orc, BoolCircuit.AdjWalk.rg, BoolCircuit.AdjWalk.rg_2, BoolCircuit.AdjWalk.rg_3, BoolCircuit.AdjWalk.upd_3, BoolCircuit.AdjWalk.walk, Complexity.LogProg.astep}
\difficulty{5}
Let $C$ be a family of directed acyclic circuits, let $lm$ be a Boolean, let $n,i,B,V,U$
be natural numbers, and let $o$ be a Boolean oracle on the five query indices and bit
strings. Suppose $V\text{,}U\le B$. For every $d,W,k$ with $W+d=U$, $k\le i$, and
$k+d+1\le B$, start the walk on the encoding of $1^n$ paired with the binary expansion
of $i$, in state $uL$, with register values $(k,V,U,W)$ and no answer recorded. If $k\le
i<k+d+1$, the walk halts with answer $lm$ or the $(i-k)$-th bit of the word consisting
of $d$ ones followed by zero, and every earlier configuration has all register values at
most $B$. If $k+d+1\le i$, the walk reaches a configuration with state $aN$, register
values $(k+d+1,V,U,U)$, and no answer recorded; every configuration before reaching it
has all register values at most $B$.
\end{lemma}

\begin{lemma}[Characterization of circuit edges]
\lean{BoolCircuit.AdjWalk.edge_iff}
\label{BoolCircuit.AdjWalk.edge_iff}
\leanok
\uses{BoolCircuit.DAGCircuit.Edge, BoolCircuit.DAGGate}
\difficulty{1}
Let $C$ be a directed acyclic circuit on $n$ input vertices, and let $g$ be a gate such
that the gate at index $V-n$ in the circuit's gate list is $g$. If $n\le V$, then $U$ is
an argument of $g$ if and only if there is an edge from vertex $U$ to vertex $V$ in the
circuit.
\end{lemma}

\begin{lemma}[Bounded traversal of gate arguments]
\lean{BoolCircuit.AdjWalk.aLoop}
\label{BoolCircuit.AdjWalk.aLoop}
\leanok
\uses{BoolCircuit.AdjWalk.Em, BoolCircuit.AdjWalk.Em.append, BoolCircuit.AdjWalk.Em.cont, BoolCircuit.AdjWalk.Em.emit, BoolCircuit.AdjWalk.Em.mono, BoolCircuit.AdjWalk.Em.step, BoolCircuit.AdjWalk.Site.bit, BoolCircuit.AdjWalk.Site.next, BoolCircuit.AdjWalk.astep_call₂, BoolCircuit.AdjWalk.edge_iff, BoolCircuit.AdjWalk.gd_rg, BoolCircuit.AdjWalk.orc, BoolCircuit.AdjWalk.orc_edge, BoolCircuit.AdjWalk.rg, BoolCircuit.AdjWalk.rg_1, BoolCircuit.AdjWalk.rg_2, BoolCircuit.AdjWalk.uLoop, BoolCircuit.AdjWalk.upd_1, BoolCircuit.AdjWalk.upd_2, BoolCircuit.AdjWalk.upd_3, BoolCircuit.AdjWalk.walk, BoolCircuit.DAGCircuit.Edge, BoolCircuit.DAGGate, BoolCircuit.encodeList, BoolCircuit.encodeNat, Complexity.LogProg.AReach.refl, Complexity.LogProg.AReach.step, Complexity.LogProg.astep}
\difficulty{6}
Let $C$ be a family of finite acyclic Boolean circuits, fix natural numbers $n,i,B$, a
Boolean parameter $\ell$, and a gate $g$ of the circuit for input length $n$. Suppose
the gate at index $V-n$ is $g$, $n\le V$, and $V+1\le B$. For every choice of natural
numbers $d,U,k,W,K$ such that $U+d=V$, $k\le i$, $W\le B$, $K=k+|E|$, and $K\le B$,
where $E$ is the encoding of the indices in $[U,U+d)$ that occur among the arguments of
$g$, the emission property holds for $E$ starting from control state $aL$, registers
$(k,V,U,W)$, and no answer. Its postcondition is that the resulting configuration has
control state $gS$, registers $(K,V+1,V,W')$, and no answer for some $W'\le B$.
\end{lemma}

\begin{lemma}[Filtering the sorted gate arguments]
\lean{BoolCircuit.AdjWalk.args_eq_filter}
\label{BoolCircuit.AdjWalk.args_eq_filter}
\leanok
\uses{BoolCircuit.DAGGate}
\difficulty{4}
Let $g$ be a gate whose argument list is strictly increasing, and suppose every argument
is less than $V$. Then filtering the list $0,1,\ldots,V-1$ to retain exactly the entries
that occur as arguments of $g$ returns the argument list of $g$.
\end{lemma}

\begin{lemma}[Vertex type of a gate]
\lean{BoolCircuit.AdjWalk.vtype_gate}
\label{BoolCircuit.AdjWalk.vtype_gate}
\leanok
\uses{BoolCircuit.DAGCircuit.size, BoolCircuit.DAGCircuit.vtype, BoolCircuit.DAGGate, BoolCircuit.GateKind}
\difficulty{3}
Let $C$ be a family of directed acyclic circuits, let $n,V$ be natural numbers, and let
$g$ be a gate such that gate number $V-n$ in the circuit for input length $n$ is $g$,
with $n\le V$. For any optional gate kind $t$, the vertex $V$ lies below the circuit
size and has type $t$ if and only if the gate kind of $g$, as an optional value, equals
$t$.
\end{lemma}

\begin{lemma}[Encoded Gate Kind Walk]
\lean{BoolCircuit.AdjWalk.kindEm}
\label{BoolCircuit.AdjWalk.kindEm}
\leanok
\uses{BoolCircuit.AdjWalk.Em, BoolCircuit.AdjWalk.Em.append, BoolCircuit.AdjWalk.Em.emit, BoolCircuit.AdjWalk.Em.mono, BoolCircuit.AdjWalk.Em.step, BoolCircuit.AdjWalk.Site.bit, BoolCircuit.AdjWalk.astep_call₁, BoolCircuit.AdjWalk.gd_rg, BoolCircuit.AdjWalk.orc, BoolCircuit.AdjWalk.orc_type, BoolCircuit.AdjWalk.rg, BoolCircuit.AdjWalk.rg_1, BoolCircuit.AdjWalk.vtype_gate, BoolCircuit.DAGCircuit.size, BoolCircuit.DAGCircuit.vtype, BoolCircuit.DAGGate, BoolCircuit.GateKind.encode}
\difficulty{5}
Let $C$ be a family of Boolean circuits, and fix Boolean $\ell_m$ and natural numbers
$n,i,B$. Suppose a gate $g$ occurs at index $V-n$ in the circuit for input length $n$,
with $n\le V\le B$. If $k\le i$, $k+2\le B$, and $U,W\le B$, then the walk initialized
in state $kA$ with registers $(k,V,U,W)$ and no result has the following behavior on the
input encoding $n$ and $i$: for $k\le i<k+2$, it halts with answer $\ell_m$ or the
corresponding bit of the two-bit encoding of $g$'s kind; for $k+2\le i$, it reaches the
specified postcondition, with every register bounded by $B$ throughout the run before
reaching it.
\end{lemma}

\begin{lemma}[Circuit gate-processing loop]
\lean{BoolCircuit.AdjWalk.gLoop}
\label{BoolCircuit.AdjWalk.gLoop}
\leanok
\uses{BoolCircuit.AdjWalk.Em, BoolCircuit.AdjWalk.Em.append, BoolCircuit.AdjWalk.Em.emit, BoolCircuit.AdjWalk.Em.mono, BoolCircuit.AdjWalk.Em.step, BoolCircuit.AdjWalk.Site.bit, BoolCircuit.AdjWalk.Site.next, BoolCircuit.AdjWalk.aLoop, BoolCircuit.AdjWalk.args_eq_filter, BoolCircuit.AdjWalk.astep_call₁, BoolCircuit.AdjWalk.gd_rg, BoolCircuit.AdjWalk.kindEm, BoolCircuit.AdjWalk.orc, BoolCircuit.AdjWalk.orc_size, BoolCircuit.AdjWalk.rg, BoolCircuit.AdjWalk.rg_1, BoolCircuit.AdjWalk.upd_2, BoolCircuit.AdjWalk.walk, BoolCircuit.DAGCircuit.IsCanonical, BoolCircuit.DAGCircuit.size, BoolCircuit.DAGGate.encode, BoolCircuit.GateKind.encode, BoolCircuit.encodeList, BoolCircuit.encodeNat, Complexity.LogProg.astep}
\difficulty{6}
Let $C$ be a family of directed acyclic circuits, and let $n,i,B$ be natural numbers and
Let $\mathit{lm}$ be a Boolean parameter. Assume the circuit $C_n$ has its gate argument
lists sorted in strictly increasing order and satisfies
$\text{output}(C_n)+1=\text{size}(C_n)$, with $\text{size}(C_n)+1\le B$. For natural
numbers $d,V,k,U,W,K$ such that $V+d=\text{size}(C_n)$, $n\le V$, $k\le i$, $U,W\le B$,
and $K=k+|S|\le B$, where $S$ is the encoding of the gates of $C_n$ from index $V-n$
onward, start the gate-processing machine in state $gS$, with registers $(k,V,U,W)$, no
recorded answer, and input encoding $1^n$ paired with the binary representation of $i$.
If $k\le i<k+|S|$, the machine halts with answer $\text{lm}\lor S_{i-k}$; if $k+|S|\le
i$, it reaches a configuration of the form $(oS,K,\text{size}(C_n),U',W',\text{no
answer})$ for some $U',W'\le B$. At every earlier step in either case, all four register
values are at most $B$.
\end{lemma}

\begin{lemma}[Input vertices have no gate type]
\lean{BoolCircuit.AdjWalk.vtype_none}
\label{BoolCircuit.AdjWalk.vtype_none}
\leanok
\uses{BoolCircuit.DAGCircuit.size, BoolCircuit.DAGCircuit.vtype}
\difficulty{4}
Let $C$ be a family of directed acyclic circuits, and let $n,V$ be natural numbers with
$V\le n$. The vertex $V$ lies below the size of the circuit for input length $n$ and has
no gate type if and only if $V<n$.
\end{lemma}

\begin{lemma}[Bounded Input Scan Loop]
\lean{BoolCircuit.AdjWalk.inLoop}
\label{BoolCircuit.AdjWalk.inLoop}
\leanok
\uses{BoolCircuit.AdjWalk.Em, BoolCircuit.AdjWalk.Em.append, BoolCircuit.AdjWalk.Em.emit, BoolCircuit.AdjWalk.Em.mono, BoolCircuit.AdjWalk.Em.step, BoolCircuit.AdjWalk.Site.bit, BoolCircuit.AdjWalk.Site.next, BoolCircuit.AdjWalk.astep_call₁, BoolCircuit.AdjWalk.gd_rg, BoolCircuit.AdjWalk.orc, BoolCircuit.AdjWalk.orc_type, BoolCircuit.AdjWalk.rg, BoolCircuit.AdjWalk.rg_1, BoolCircuit.AdjWalk.upd_1, BoolCircuit.AdjWalk.vtype_none, BoolCircuit.AdjWalk.walk, Complexity.LogProg.astep}
\difficulty{5}
Let $C$ be a family of Boolean circuits, and fix Boolean $\ell_m$ and natural numbers
$n,i,B$. Suppose $n+1\le B$. For all natural numbers $d,V,k,U,W$ with $V+d=n$, $k\le i$,
$U,W\le B$, and $k+d+1\le B$, the walk initialized in state $nI$ with registers
$(k,V,U,W)$ and no result has the following behavior on the input encoding $n$ and $i$:
if $k\le i<k+d+1$, it halts with answer $\ell_m$ or the $(i-k)$-th bit of the word
consisting of $d$ ones followed by zero; if $k+d+1\le i$, it reaches a configuration
satisfying the specified postcondition, with all registers bounded by $B$ throughout the
run before reaching it.
\end{lemma}

\begin{lemma}[Output Scan Loop]
\lean{BoolCircuit.AdjWalk.oLoop}
\label{BoolCircuit.AdjWalk.oLoop}
\leanok
\uses{BoolCircuit.AdjWalk.Em, BoolCircuit.AdjWalk.Em.append, BoolCircuit.AdjWalk.Em.emit, BoolCircuit.AdjWalk.Em.mono, BoolCircuit.AdjWalk.Em.step, BoolCircuit.AdjWalk.Site.bit, BoolCircuit.AdjWalk.Site.next, BoolCircuit.AdjWalk.gd_rg, BoolCircuit.AdjWalk.orc, BoolCircuit.AdjWalk.rg, BoolCircuit.AdjWalk.rg_1, BoolCircuit.AdjWalk.rg_3, BoolCircuit.AdjWalk.upd_3, BoolCircuit.AdjWalk.walk, Complexity.LogProg.astep}
\difficulty{4}
Let $C$ be a family of Boolean circuits, and fix Boolean $\ell_m$ and natural numbers
$n,i,B$. For every $V\le B$, if $W+d=V$, $k\le i$, $U\le B$, and $k+d+1\le B$, then the
walk initialized in state $oL$ with registers $(k,V,U,W)$ and no result has the
following behavior on the input encoding $n$ and $i$: if $k\le i<k+d+1$, it halts with
answer $\ell_m$ or the $(i-k)$-th bit of the word consisting of $d$ ones followed by
zero; if $k+d+1\le i$, it reaches the specified postcondition, with registers
$(k+d+1,V,U,V)$ and state $done$, while all registers are bounded by $B$ throughout the
run before reaching it.
\end{lemma}

\begin{theorem}[Termination of the circuit walk]
\lean{BoolCircuit.AdjWalk.walk_halts}
\label{BoolCircuit.AdjWalk.walk_halts}
\leanok
\uses{BoolCircuit.AdjWalk.Em, BoolCircuit.AdjWalk.Em.append, BoolCircuit.AdjWalk.Em.mono, BoolCircuit.AdjWalk.Em.step, BoolCircuit.AdjWalk.Gd, BoolCircuit.AdjWalk.X, BoolCircuit.AdjWalk.gLoop, BoolCircuit.AdjWalk.gd_rg, BoolCircuit.AdjWalk.inLoop, BoolCircuit.AdjWalk.oLoop, BoolCircuit.AdjWalk.orc, BoolCircuit.AdjWalk.rg, BoolCircuit.AdjWalk.rg_3, BoolCircuit.AdjWalk.rg_init, BoolCircuit.AdjWalk.upd_3, BoolCircuit.AdjWalk.walk, BoolCircuit.DAGCircuit.IsCanonical, BoolCircuit.DAGCircuit.encode, BoolCircuit.DAGCircuit.encode_eq, BoolCircuit.DAGCircuit.size, BoolCircuit.DAGCircuitFamily, BoolCircuit.DAGGate.encode, BoolCircuit.encodeList, BoolCircuit.length_encodeList_ge, Complexity.LogProg.AHalt, Complexity.LogProg.AHalt.ret, Complexity.LogProg.AHalt.step, Complexity.LogProg.AReach.halt, Complexity.LogProg.ValidPlain, Complexity.LogProg.astep, Complexity.LogProg.canon_bits}
\difficulty{5}
Let $C$ be a family of directed acyclic circuits, and suppose the circuit for input
length $n$ is canonical: each gate's argument list is strictly increasing, and its
output vertex is the last vertex. For any Boolean parameter $\ell$ and natural number
$i$, the walk on the paired encoding of the unary word $1^n$ and the binary
representation of $i$, started in its start state with all four registers zero and no
answer yet, eventually halts with answer $\ell\lor b$ if $i$ is less than the circuit's
encoding length, where $b$ is bit $i$ of that encoding (taking a missing bit to be
false); otherwise it halts with answer false. At every time strictly before it halts,
each register is at most the length of the circuit encoding.
\end{theorem}
